\section{Discussion}\label{sec:discussion}

\subsection{What the transfer ratio shows}
On single-hop NQ, a retrieval gain of a dense encoder over a small encoder reached the answer at a rate of \ev{T.lnq.R4} to \ev{T.lnq.R3} EM points per Hit@5 point, with both readers. Three observations within our setup are consistent with this. (i) The readers fell well short of their gold-passage bound even with good retrieval, so part of the loss lies in reading, not retrieval. (ii) EM given a miss stayed far above closed-book EM. The NQ relevance labels are sparse, so many ``misses'' likely still contained an answer-bearing passage that was not labelled gold, or the reader knew the answer. Hit@5 may therefore overstate the difference between encoders that matters for the reader. (iii) EM given a hit fell slightly as encoders improved, which suggests that the questions that a better encoder newly solves at the retrieval stage are harder to answer. This last point is a hypothesis; we did not test which questions changed status.

On multi-hop HotpotQA, Hit@5 saturated while Both@5 did not. In the baseline-conditioned decomposition, between \ev{DS.hp.min} and \ev{DS.hp.max} of the transfer came from higher EM given a hit (\ev{V4.sym.hp.min}--\ev{V4.sym.hp.max} in the symmetric form), and the EM ranking of retrievers followed Both@5 exactly. Both are descriptions of the data, not causal attributions. A reasonable reading is that, for bridge questions, the second supporting passage carries the answer, and Hit@5 does not count it. A post hoc check supports this reading: with Both@5 as the denominator, the HotpotQA ratios fell to \ev{V4.both.hp.min}--\ev{V4.both.hp.max}, much closer to the NQ range. The large Hit@5-based ratios on HotpotQA therefore largely reflect the denominator: the retrieval gain that matters is larger than Hit@5 shows. The comparison questions, where both entities are named in the question, were close to their gold bound for every dense encoder. Teams that build multi-hop systems may therefore prefer to measure retrieval with a metric that requires all needed evidence.

Measured with Hit@5, the two datasets thus support different conclusions about the same encoders. A team that judges encoders by single-hop retrieval benchmarks may overrate the answer value of the strongest ones. A team with multi-hop questions may underrate it if it uses Hit@k.

\subsection{Comparison with earlier work}
Our NQ results agree with reports that retrieval gains do not carry over fully to answers \citep{salemi2024erag,mathur2026policy,ghosh2026evidence} and with the falling recall conversion rate reported by \citet{peng2026neocorrag}. They do not contradict the strong correlations reported by \citet{samuel2026beyond}: the EM ranking of retrievers followed their Hit@5 ranking closely in every cell (Spearman $\rho$ between \ev{RHO.lnq} and \ev{RHO.qnq}). Correlation concerns the order of systems; the transfer ratio concerns the size of the gain. A metric can rank encoders correctly and still overstate how much answers improve.

The context experiment agrees with earlier findings on distraction. Top-ranked non-gold passages lowered EM compared with random passages \citep{cuconasu2024power,amiraz2025distracting,repantis2026competition}, and adding random passages to the gold passage changed EM by no detectable amount. Gold position showed no detectable effect among five passages, in line with \citet{cuconasu2025positional}; with five short passages our setting is far from the long contexts in which position effects were originally shown \citep{liu2024lost}. The near-identical retrieval results of exact and approximate search agree with \citet{leto2024toward}, who also measured answers; we measured retrieval only.

\subsection{Practical implications}
Within our setup, the choice of encoder mattered less for answers than for retrieval, while its cost differed by more than an order of magnitude. For single-hop questions, a base-size encoder such as e5-base-v2 or bge-base gave answers that no larger encoder improved detectably, while encoding \ev{E7.ratio.R5R4} times faster than bge-large. For multi-hop questions, bge-large gave a small but detectable gain. Qwen3-Embedding-0.6B was the most expensive encoder and gave lower EM than e5-base-v2 on HotpotQA; its instruction and length settings may not suit this task, and we did not tune them. Retrieving ten passages instead of five raised EM for most retrievers and narrowed their differences, which is another lever when the reader's context allows it, although it nearly doubles the prompt the reader must process. A cross-encoder helped the weak BM25 retrieval stage but not a strong dense one. Approximate search returned nearly the same top five passages as exact search at a small fraction of the latency; we did not measure its effect on answers.

\subsection{Alternative explanations}
Several features of the design could produce low NQ transfer without a property of the encoders. EM undercounts correct answers \citep{kamalloo2023}; the F1 and containment ratios were higher but showed the same contrast between datasets. The sub-corpus inflates dense retrieval scores, but the full-corpus check kept the order of four retrievers, and the inflation was similar for the three dense encoders tested (\ev{E9.dropabs.R2} to \ev{E9.dropabs.R1} in nDCG@10). The construction favours the dense encoders relative to BM25, so the comparisons that involve BM25 depend most on it. The readers were 4-bit quantised 7--8B models with one prompt; stronger readers may convert hits into answers at a different rate. The baseline encoder also matters: $T(R)$ against MiniLM-L6 measures gains above a fairly strong small encoder. Against BM25 the NQ ratios were higher (\ev{V4.b0.nq.min}--\ev{V4.b0.nq.max}), and the contrast between datasets held for four of the five dense encoders but reversed for Qwen3-Embedding-0.6B (Table~\ref{tab:altbase}). The ratio therefore describes a pair of retrievers, not a single encoder, and the dataset contrast is clearest against the small dense baseline.
